% Français 9115, batch 28 — Gallica views 541–560.
% Pass: Opus 5.5 (claude-opus-5-5), 2026-09-22 — first pass, unchecked against the views by a human.
%
% What the twenty views hold. Three kinds of paper, all hers as far as the
% leaves show, none of it number theory.
%
% (1) v. 541–548: eight pages of a paginated memoir in French, in her
% regular fair hand, corrected between the lines, paginated (131)–(137) at
% the head, centre. The leaves are written on both sides; the ink series
% stands on the odd-paginated side. It is the experimental part of a
% memoir on vibrating elastic plates, and plainly the same run as the
% prize-memoir draft that batch 22 left at page (34), v. 440, and it
% continues batch 27 without a break: pages (111)–(130) at v. 521–540,
% same layout, ink 256–265 (read in that file once it appeared, after
% this pass had read its own views). The wax-loaded plate experiment
% begins on v. 539–540, in batch 27.
%   v. 541  (131) square plate loaded with wax (1 gros 2 grains, an eighth
%           of the plate's weight of 1 once 18 grains) on one side, then on
%           two contiguous sides, then on two opposite sides; the tones
%           heard (« entre re et re# 1 », « exactement re 1 ») and the
%           displacement of the nodal lines, measured in millimetres.
%   v. 542  (132) « explication »: each dose of wax is equivalent to
%           106/8 ≈ 13 mm added to the side; one dose turns the square of
%           106 mm into a rectangle 106 : 119 ≈ 8 : 9, the tone
%           proportional to (9²+8²)/64 = 145/64.
%   v. 543  (132, repeated) the ratio of tones 145/128 hardly differs from
%           9/8, a major tone, as heard; the nodal line at 119/2 mm, 58 and
%           45 mm observed; two doses make a square of 119 mm, 9²/8², two
%           major tones.
%   v. 544  (133) the same for two opposite sides: a rectangle 10 : 8, tone
%           41/16, ratio 41/32, a little over two major tones; the nodal
%           line does not move.
%   v. 545  (134?) the remarkable fact: a small part between the loaded
%           side and the nearest nodal line gives the same tone as a much
%           larger part between two lines; a glass lamina 20 cm × 23 mm,
%           4 gros 12 grains, loaded with 1 gros 3 grains at one end: 12 mm
%           and 121 mm give the same sound.
%   v. 546  (135) « une de ces portions est dix fois plus étendue que
%           l'autre »; in general a portion as small as one likes can be
%           made to give the tone of a given portion by adding non-elastic
%           matter; the physical explanation begun.
%   v. 547  (136) the explanation: the added weight stands for the
%           prolongation of the plate; the analysis shows the wax may be
%           applied at any distance from the origin, the terms that must
%           vanish at the limits being multiplied by a factor depending on
%           the rigidity.
%   v. 548  (137) only rectangular plates have been treated; a new section,
%           « § 7. Recherche d'une intégrale particulière immédiatement
%           applicable au cas d'une plaque elastique circulaire », N.o 21
%           (read with doubt), polar coordinates r and φ in the plane of the
%           plate. The page stops mid-sentence on « on pourra toujours
%           faire »; its sequel is not in this batch (see (3)).
% (2) v. 549: a single draft sheet, a different paper, a fast hand in two
%   inks, heavily struck, with trial words in the left margin and the
%   Bibliothèque's stamp: not elasticity but the motion of heat in a
%   cylinder — an equation dv/dr + (h/k)v = 0 labelled (D), « les
%   equations (B) et (D) », « l'analyse employée N.o 19 » and « N.o 14 »,
%   the terms K/CD (d²v/dr² + (1/r) dv/dr) and K/CDrr dv/dφ of an equation
%   (C). One isolated line at its head, z = … {e^{ωx/a} − e^{(a−ω)x/a} + (1,
%   in a paler ink. Whose argument this is (her own, or notes on someone
%   else's heat memoir) the sheet does not say.
% (3) v. 553–560: calculation leaves for the circular plate of § 7, with
%   no prose to speak of. v. 553: two small slips on one mount — the upper,
%   inked 273, the partial derivatives in r, φ, δ; the lower, inked 272,
%   the limit equation of the round plate satisfied by sin φ = 0, and
%   d²z/dxdy = 0 giving z = F(x) + f(y). v. 555: a slip inked 274,
%   « cas particuliers de l'intégrale (E) », the particular integral in
%   sin(πnx/A) sin(πmy/A) and its derivatives. v. 557 (page (138), struck,
%   ink 275) and v. 560 (its back, page (139)): the fair transformation of
%   the plate equation into polar coordinates — dz/dx, dz/dy,
%   d⁴z/dx⁴ + 2d⁴z/dx²dy² + d⁴z/dy⁴ in r and φ — then, with bτ² = e⁴/a⁴,
%   « l'equation différentielle (B) (V. N.o 2) », z = u sin(βφ + B), the
%   equation (B') in u, and u = r^β s expanded to d³u/dr³. The page
%   numbers (138)–(139) continue (137) at v. 548, so these two pages are the
%   sequel of § 7; the slips between are working papers bound among them.
%
% Views not transcribed. v. 550: the back of the sheet of v. 549, blank,
% the draft showing through reversed. v. 551–552: a leaf inked 271, blank
% on both sides. v. 554, 556, 558: blank backs of the mounting sheets,
% on which the upper slip of v. 553 shows through (checked by mirroring
% v. 556 and v. 558). v. 559: the page of v. 557 photographed a second
% time (same text, same « (138) », same 275). The versos of v. 541–547
% that are written are transcribed as v. 542, 544, 546, 548; none is blank.
%
% Seams. View 541 opens a new paragraph on a new page, (131); the page
% before, (130) at v. 540 (not read, glanced at), ends a paragraph on
% « l'extremité opposée au côté chargé. », so nothing is cut at 540/541;
% batch 27, now written, reads the same seam the same way.
% View 548 stops mid-sentence (« on pourra toujours faire ») and its
% sequel is not on the next leaves of this batch. View 560 ends on the
% expansion of d³u/dr³ at the foot of a page; view 561 is a new slip inked
% 276, « On a … », which may carry the calculation on (d⁴z/dr⁴ terms) —
% not settled from the thumbnail.
%
% The numbers on the leaves, and why no \folio{} is written. (1) Her
% pagination, in parentheses at the head, centre: (131) v. 541, (132)
% v. 542, (132) again v. 543, (133) v. 544, (13?) v. 545 — the last figure
% rewritten, most likely 4 over 2 — (135) v. 546, (136) v. 547, (137)
% v. 548, (138) v. 557 (struck through with a horizontal stroke, and
% again on its repeat v. 559), (139) v. 560. Her 3 and 9 are close: the
% middle figure is read 3 throughout because (138), (139) follow (137).
% Two consecutive pages read (132); recorded as read. (2) The ink series
% at the top right, in the larger rounder hand: 266 (v. 541), 267 (543),
% 268 (545), 269 (547), 270 (549, the heat draft), 271 (551, a blank
% leaf), 273 (upper slip, v. 553), 272 (lower slip, v. 553), 274 (555),
% 275 (557 and its repeat 559); never on the backs (542, 544, 546, 548,
% 560). One number per physical leaf, slips and a blank leaf included,
% as in batch 22 — and as there penned, not pencilled, so no \folio{}.
% The two slips on one mount are numbered out of order from top to
% bottom (273 above 272). It continues 216 at v. 441 at one leaf per two
% views. (3) Her paragraph numbers: « § 7 » and N.o 21 (doubtful, perhaps
% 31) at v. 548; references to N.o 2 (v. 560), N.o 14 and N.o 19 (v. 549,
% in the heat draft). Every number above was read on its view; none is
% inferred. None of these leaves can be Laubenbacher and Pengelley's
% « Manuscript C » (ff. 348r–349r): the subject is plates and heat, and
% the ink series here runs 266–275.
%
% Notation. Hers, kept. Tones by note names with a figure after them (« re
% 1 », « mi 1 et fa 1 », « sol 1 »): the number of the octave, counted
% from the octave of the tuning-fork la as she defines it at v. 539
% (batch 27); the sharp is set $\sharp$. Weights in « once », « gros », « grains »;
% lengths in millimetres. In the calculations, round brackets round a
% partial derivative are hers; « sin », « cos » (with or without a stop)
% are set \text{sin}, \text{cos}; φ and δ are her angles; the π of v. 555
% is a small n with a bar over it, followed by a stop (« π.nx »). The
% capital read D in the heat draft (v. 549) is a dome on a bar and might
% be an Ω. Words she underlines in the prose are set in \emph{}; letters
% underlined inside the mathematics keep \underline{}. Plural endings
% cancelled with her small cross (« parallelle », « nodale »,
% « transversale », « son », « être », « soumis », « matière »,
% « memoire », « indispensable ») are given in the singular with a note.
% Spelling is hers and is not flagged: « parallelle », « précédant »,
% « équivalans », « equivalante », « foible », « fesant », « quarrés »,
% « serappellant », « parceque », « tems », « differens », « dela »,
% « aulieu », « ceque », « cequi », « desorte », « apeuprès », « il ya »,
% « c'est-àdire », « le quel », « Apresent ». Where the leaf and the
% calculation disagree the leaf stands, with a note (v. 553, 555, 557).
%
% What could not be read, and what is offered with doubt. \ill{} stands
% fourteen times: in the heat draft (v. 549) a small abbreviated word
% added above « même », a struck word in « de la … dela chaleur », two
% struck groups after « employée dans cette », and a struck word before
% « (C) » in the margin; the blotted exponent of r (v. 553); on the slip
% of v. 555 a blotted sign between the two brackets of z and three blotted
% signs in the second derivative (named a fifth time in the slip's note);
% in v. 560 three blackened groups of the struck first attempt at the
% equation. \uncertain{} stands for « six » (v. 542, struck), « environ »
% (v. 542, blotted), « 1 » after sol and « repose » (v. 543, struck), « à »
% and « la » (v. 544), « 2 » (v. 545), « paroît » (its ending) and
% « lagene » (v. 546), « et », « equivalante », « satisfera » (v. 547),
% N.o « 21 » (v. 548), « +dv/dr », « deux », N.o « 19 » and « 14 »
% (v. 549), ∫∫ (v. 553), « n = −1 » (v. 555), the last denominator
% dφ⁴ (v. 557). Nothing was guessed to fill a gap.
%
% Nothing here was checked against any published edition of Sophie
% Germain's papers or of any memoir on plates or on heat.
\documentclass[11pt,a4paper]{article}
% The preamble shared by every transcription of the mathematicians' manuscripts in Gallica.
%
% It rests on one idea: **uncertainty compounds, it does not resolve.** A
% transcription that smooths over a doubtful word has destroyed the only thing
% it was made for — a reader must be able to tell, at every point, what was on
% the page and what was supplied. The apparatus macros below are therefore the
% critical apparatus, not decoration.
%
% The same commands are recognised by `scripts/render.mjs`, which produces the
% site's reading view, and by `scripts/tei.mjs`. A macro added here must be
% added there too, or rendering fails loudly — which is the wanted behaviour.
%
% Comments here are English, like the rest of the repository. The documents
% this preamble sets are French, and that is deliberate: see the skills.

\ProvidesPackage{germain}[2026/09/15 Transcriptions of mathematicians' manuscripts in Gallica]

\usepackage{amsmath,amssymb,amsthm}
\usepackage{mathrsfs}
% Kept from the parent preamble so the renderer's diagram support stays
% available; these manuscripts draw no commutative diagrams, and nothing here uses it.
\usepackage{tikz-cd}
\usepackage[english,french]{babel}
\usepackage{fontspec}

% --- Greek ------------------------------------------------------------------
% The text face stays fontspec's default, Latin Modern, which has no Greek:
% XeTeX drops every glyph it lacks with a line in the log and nothing on the
% page, and the Greek words the manuscripts quote vanished from the PDFs. So
% the Greek and Greek Extended blocks get a XeTeX character class of their own,
% and entering or leaving a run of it switches the family to CMU Serif — the
% same Computer Modern design, with polytonic Greek — and back. Only the
% family changes: a Greek word in \textit or \textbf keeps its shape and series.
% The transitions to and from every other class carry the tokens that class
% has towards an ordinary letter, so French spacing before « ; » or inside
% « » still holds after a Greek word. No grouping, so no brace or \endgroup
% can come out unbalanced; maths is left alone, since XeTeX inserts nothing
% there. Kept identical in `exercices.sty`.
\makeatletter
\newfontfamily\germainGreekfont{cmun}[
  Extension=.otf, NFSSFamily=germaingreek,
  UprightFont=*rm, ItalicFont=*ti, SlantedFont=*sl,
  BoldFont=*bx, BoldItalicFont=*bi, BoldSlantedFont=*bl]
\newXeTeXintercharclass\germain@greek
\def\germain@greekrange#1#2{%
  \@tempcnta=#1\relax
  \loop
    \XeTeXcharclass\@tempcnta=\germain@greek
    \ifnum\@tempcnta<#2\advance\@tempcnta\@ne\repeat}
\germain@greekrange{"0370}{"03FF}
\germain@greekrange{"1F00}{"1FFF}
\def\germain@greekprev{\rmdefault}
\def\germain@greekin{%
  \edef\germain@greekprev{\f@family}\fontfamily{germaingreek}\selectfont}
\def\germain@greekout{\fontfamily{\germain@greekprev}\selectfont}
\def\germain@greekjoin{%
  \XeTeXinterchartokenstate=\@ne
  \@tempcnta=\z@
  \loop
    \ifnum\@tempcnta=\germain@greek\else
      \XeTeXinterchartoks\germain@greek\@tempcnta=\expandafter{%
        \expandafter\germain@greekout\the\XeTeXinterchartoks\z@\@tempcnta}%
      \XeTeXinterchartoks\@tempcnta\germain@greek=\expandafter{%
        \the\XeTeXinterchartoks\@tempcnta\z@\germain@greekin}%
    \fi
    \ifnum\@tempcnta<4095\advance\@tempcnta\@ne\repeat}
% babel-french sets its own classes, and saves and restores the interchar
% state, each time a language is selected: join again after it does.
\addto\extrasfrench{\germain@greekjoin}
\addto\extrasenglish{\germain@greekjoin}
\AtBeginDocument{\germain@greekjoin}
\makeatother

\usepackage{geometry}
\geometry{a4paper,margin=25mm,marginparwidth=28mm}
\usepackage{marginnote}
\usepackage{xcolor}
\usepackage[normalem]{ulem}
\setlength{\emergencystretch}{3em}
\usepackage{hyperref}
\hypersetup{colorlinks=true,linkcolor=black,urlcolor=black,pdfborder={0 0 0}}

% --- Batch metadata ---------------------------------------------------------
% Taken as they stand by the rendering script, which shows them at the head of
% the reading view. They fix what the file claims to cover. The macro names are
% the parent project's, so that the scripts stay shared; \folder here means the
% volume's slug — `fr-9115` — and \pages counts Gallica views.

\newcommand{\folder}[1]{\gdef\grothFolder{#1}}
\newcommand{\batch}[1]{\gdef\grothBatch{#1}}
\newcommand{\pages}[2]{\gdef\grothFirst{#1}\gdef\grothLast{#2}}
\newcommand{\foldertitle}[1]{\gdef\grothFoldertitle{#1}}

% The holder's own shelfmark, printed as they write it: « Français 9115 ».
\newcommand{\shelfmark}[1]{\gdef\grothShelfmark{#1}}

% The dating, copied verbatim from the holder's catalogue — never inferred
% here. « XIXe siècle » is the BnF's; a pass that dates a leaf from its content
% says so in a \note{}, as its own inference.
\newcommand{\dating}[1]{\gdef\grothDating{#1}}

% The status notice, printed in the title block and repeated as a diagonal
% watermark on every page of the PDF. The manuscripts are in the public
% domain; what this line declares is not a right but a quality — an
% unverified first pass by a machine. The skills write the line; a file
% without it gets no watermark, so leaving it out is visible in the output.
\newcommand{\watermark}[1]{\gdef\grothWatermark{#1}}

\gdef\grothFolder{?}\gdef\grothBatch{}\gdef\grothFirst{?}\gdef\grothLast{?}%
\gdef\grothFoldertitle{}\gdef\grothShelfmark{}\gdef\grothDating{}\gdef\grothWatermark{}

% --- The résumé -------------------------------------------------------------
\newenvironment{resume}
  {\small\setlength{\parskip}{0.45em}}
  {\par\medskip\hrule\medskip}

% --- Keywords ---------------------------------------------------------------
% In English, closing the modernised reading's résumé: the modern vocabulary
% under which to search for what the leaves construct. The manifest extracts
% them to tag the volume — this line is the single source of the tags.
\newcommand{\keywords}[1]{\par\smallskip\noindent
  {\footnotesize\textsc{Keywords} --- \textit{#1}}\par}

% --- The critical apparatus -------------------------------------------------

% The Gallica view number, in the margin. This is the anchor that lets a
% disagreement over a formula be settled by looking at *one* image rather than
% a volume of 750. The site uses it to turn the facsimile as the transcription
% is scrolled. A view is one scanned image: usually one side of a leaf.
\newcommand{\page}[1]{%
  \marginnote{\footnotesize\color{gray!70}\textsf{v.\,#1}}%
  \label{page:#1}\ignorespaces}

% The BnF's pencilled foliation, where the leaf carries it and a pass can read
% it: « 198r ». This is what the literature cites — « ff. 198r–208v » — and
% the one line that turns a citation into a view. Recorded, never inferred:
% a folio a pass cannot read is not written.
\newcommand{\folio}[1]{%
  \marginnote{\footnotesize\color{gray!50}\textsf{f.\,#1}}\ignorespaces}

% The modernised reading's coarser counterpart to \page{}: which views a
% section is drawn from.
\newcommand{\pagerange}[2]{%
  \marginnote{\footnotesize\color{gray!70}\textsf{v.\,#1--#2}}%
  \ignorespaces}

% Illegible: never guessed.
\newcommand{\ill}{\textcolor{red!60!black}{[\,\ldots\,]}}

% A doubtful reading: the word is given, and flagged as uncertain.
\newcommand{\uncertain}[1]{\uline{#1}}

% An editorial addition — a missing letter, an implied word.
\newcommand{\add}[1]{\textcolor{blue!50!black}{[#1]}}

% Struck out by the writer. What was crossed out is part of the document.
\newcommand{\struck}[1]{\sout{#1}}

% The transcriber's note, on the state of the page or the mathematics.
\newcommand{\note}[1]{\par\smallskip{\small\color{gray!60!black}\textit{[#1]}}\par\smallskip}

% A marginal note of hers, distinct from the body of the text.
\newcommand{\marginal}[1]{\par\smallskip\noindent\hspace{1em}%
  {\small\color{gray!60!black}#1}\par\smallskip}

% --- Title ------------------------------------------------------------------
% The title block is French, like the documents themselves.

\AddToHook{shipout/background}{%
  \ifx\grothWatermark\empty\else
    \begin{tikzpicture}[remember picture, overlay]
      \node[rotate=52, scale=3.4, align=center, text=gray!18,
            font=\sffamily\bfseries]
        at (current page.center) {\grothWatermark};
    \end{tikzpicture}%
  \fi}

\AtBeginDocument{%
  \begin{center}
    {\large\bfseries Manuscrits de mathématiciens (Gallica) ---
      \ifx\grothShelfmark\empty\grothFolder\else\grothShelfmark\fi}\\[2pt]
    {\itshape\grothFoldertitle}\\[4pt]
    {\small vues \grothFirst--\grothLast
      \ifx\grothBatch\empty\else\ (lot \grothBatch)\fi}\\[2pt]
    \ifx\grothDating\empty\else
      {\small\color{gray!60!black}Datation du catalogue : \grothDating}\\[2pt]
    \fi
    \ifx\grothWatermark\empty\else
      {\small\bfseries\color{red!45!gray}\grothWatermark}\\[2pt]
    \fi
    {\footnotesize\color{gray!60!black}Transcription établie d'après les images IIIF de Gallica.
     Source gallica.bnf.fr / Bibliothèque nationale de France.\\
     Les lectures incertaines sont soulignées, les illisibles notées [\,\ldots\,],
     les ajouts entre crochets ; \textsf{v.} numérote les vues de Gallica,
     \textsf{f.} la foliotation au crayon de la BnF.}
  \end{center}
  \vspace{1em}}

\endinput


\folder{fr-9115}
\shelfmark{Français 9115}
\batch{28}
\pages{541}{560}
\dating{1801-1900}
\watermark{Lecture automatique\\non vérifiée}
\foldertitle{Papiers de Sophie GERMAIN. — Recueil de dissertations et problèmes mathématiques et physiques.}

\begin{document}

\note{Numérotation des feuillets. Sa pagination, entre parenthèses en tête et au milieu : (131) vue 541, (132) vue 542, (132) de nouveau vue 543, (133) vue 544, (13?) vue 545, le dernier chiffre récrit, sans doute 4 sur 2, (135) vue 546, (136) vue 547, (137) vue 548, (138) vue 557, barré d'un trait, (139) vue 560. En haut à droite, à l'encre, d'une écriture plus ronde, la série qui court dans tout le volume : 266 (vue 541), 267 (543), 268 (545), 269 (547), 270 (549), 271 (551, feuillet blanc), 273 et 272 (les deux petits feuillets de la vue 553, de haut en bas), 274 (555), 275 (557, repris à la vue 559). Cette série numérote des feuillets, petits feuillets et feuillet blanc compris ; elle occupe la place de la foliotation de la Bibliothèque, mais elle est tracée à la plume et non au crayon, et aucun « f. » n'est donc porté. Tous ces nombres ont été lus sur la vue ; aucun n'est déduit. Ne sont pas transcrites : la vue 550, dos blanc de la feuille de la vue 549 ; les vues 551 et 552, feuillet 271 blanc des deux côtés ; les vues 554, 556 et 558, dos blancs des feuilles de montage, où le petit feuillet 273 se lit par transparence ; la vue 559, seconde prise de vue de la page de la vue 557.}

\page{541}

En laissant toujours un premier côté chargé de cire j'en ai ajouté
une egale quantité, savoir 1 gros 2 grains, que j'ai collée dans
toute l'étendue d'un second côté contigu au premier.

En tirant toujours le son de l'angle libre, il s'est trouvé
entre re et re$^{\sharp}$ 1.
\note{Après le nom de la note, un trait vertical appuyé, lu « 1 », comme dans « entre mi et fa 1 » à la page précédente (vue 540, hors de ce lot) ; c'est, semble-t-il, le numéro de l'octave, compté à partir de celle du la du diapason, comme elle le définit à la vue 539. Le dièse est écrit « $\sharp$ » au-dessus de la ligne.}

La ligne nodale parallelle au second côté chargé, s'est
rapprochée de ce côté, de manière que l'intersection
des deux lignes nodales s'est faite à environ 45 millimetres
de l'un et l'autre côté chargés.

En laissant encore le premier côté chargé comme
je l'ai expliqué plus haut, aulieu d'appliquer
la seconde dose de cire à un des côtés contigus
à ce \add{premier} côté, \struck{je} j'ai fait la même application
au côté opposé à ce même premier côté.

Le son étoit alors un plus grave que dans le cas
précédant, car il étoit \struck{juste} \add{exactement} re 1.
\note{« un plus grave » : ainsi sur le feuillet, sans mot entre « un » et « plus ».}

La figure nodale est restée absolument la même
que sur la plaque non chargée; seulement les
lignes nodales occupoient plus d'espace parceque
le son étoit plus foible qu'avant l'addition de
la cire.

Une précaution indispensable\note{Le pluriel « indispensables » est annulé par un petit « x » sur la dernière lettre ; on donne le singulier.} dans les experiences de rigueur
est de laver souvent la plaque avec une liqueur

\page{542}

spiritueuse, parceque à cause dela foiblesse des sons, la
moindre trace de cire \struck{sur colleroit la poussière et}
empêcheroit les figures de se former
\note{Un long trait continu barre « sur colleroit la poussière et » ; « sur » est en outre repassé d'un trait plus lourd.}

\emph{explication}

Pour se rendre compte des \struck{\uncertain{six}} résultats des experiences précédentes,
il faut remarquer que le poids dela plaque non chargée
étant 1 once 18 grains, 1 gros 2 grains que pese\struck{nt} chaque
portion de cire ajoutée, est aussi exactement qu'il
convient à l'experience, la huitième partie de ce poids,
desorte qu'en serappellant que la longueur du côté
dela même plaque est 106 millimetres, on
trouvera que chaque portion de cire ajoutée
doit équivaloir à $\frac{106}{8}$ millimetres $=$ \uncertain{environ} 13 millim.
de prolongement du côté auquel l'application
a été faite.
\note{Une tache d'encre couvre le milieu du mot lu « environ ».}

Cela posé, il est aisé de voir cequi doit arriver
dans chacun des cas que j'ai soumis à l'experience.
D'abord en fesant une seule addition, la plaque
qui étoit de 106 millimetres quarrés, est fictivement
transformée en une plaque rectangulaire dont les
côtés sont entre eux $:: 106 : 119$ à très peu près
$:: 8 : 9$. Le son doit donc être proportionnel au
nombre $\frac{9^2+8^2}{64} = \frac{145}{64}$ tandis que dans la plaque

\page{543}

non chargée le son qui accompagne la même figure
est proportionnel au nombre 2. \struck{et la difference devra} \add{Le rapport de ces deux sons}
\struck{être} \add{sera exprimé par} $\frac{145}{128}$ qui \struck{ne} diffère \struck{pas} \add{à peine} \struck{d'une quantité appréciable}
de \struck{rapport} $\frac{144}{128} = \frac{9}{8}$, c'est-àdire d'un ton majeur.
\note{« Le rapport de ces deux sons » est écrit au-dessus de « et la difference devra », barré ; « sera exprimé par » au-dessus de « être », barré ; « à peine » au-dessus de « pas », barré avec la fin de la ligne.}

\struck{Et} En effet, tandis que le son dela plaque non
chargée est un peu plus bas que sol \uncertain{1}, celui de
la plaque dont un des côtés est chargé, est entre
mi 1 et fa 1.
\note{Après « sol », une tache d'encre couvre un trait vertical, où l'on attend le même « 1 » qu'après « mi » et « fa ».}

A l'égard dela figure, \struck{qui} \struck{\uncertain{repose}} la ligne nodale
parallelle au côté chargé, doit être à $\frac{119}{2}$ \add{$=59+\frac{1}{2}$} millimetres
du côté opposé, et à $59+\frac{1}{2}-13 = 46+\frac{1}{2}$ millimetres
du côté chargé. L'experience ne donne, àla vérité,
que 58 millimetres pour une de ces distances et
45 pour l'autre, mais il faut observer qu'il ya
environ 3 millimetres d'occupés tant par la ligne
nodale que par l'application dela cire.

En fesant une seconde addition au côté contigu au
premier, la plaque, \struck{que} \add{qui} non chargée étoit de
106 millimetres quarrés, est fictivement transformée
en une plaque de 119 millimetres quarrés; le
rapport des sons de ces deux plaques est donc exprimé
par $\frac{9^2}{8^2} =$ 2 tons majeurs. en effet j'ai trouvé un

\page{544}

de ces sons plus bas que \emph{sol} 1 et l'autre entre \emph{re} 1 et
\emph{re}$^{\sharp}$ 1.

Par la même raison que dans le cas precedent,
chacune des lignes nodales doit être à $59+\frac{1}{2}$ millimetres
du côté chargé; et l'experience y est conforme en
ayant toujours égard à l'espace occupé par ces
lignes et par l'application dela cire.

Enfin lorsque l'on fait l'addition \uncertain{à} deux \struck{côté}
des côtés opposés,\note{« à » paraît corrigé sur « au » : un accent grave ajouté, la fin du mot barrée.} la plaque est fictivement augmentée
de 13 millimetres à chacune de ces extremités; desorte
qu'elle est transformée en une plaque rectangulaire
dont les côtés sont entre eux à fort peu près
$:: 10 : 8$.

Ainsi le son doit être proportionnel au nombre
$\frac{10^2+8^2}{8^2} = \frac{41}{16}$, et \struck{pour} \uncertain{la} \struck{rapport} difference de
ce son à celui de la plaque non chargée, est
exprimée par $\frac{41}{32}$, c'est-àdire qu'elle est un
peu $>$ que 2 tons majeurs. L'experience donne
en effet de \emph{re} 1 à \emph{sol} 1 – aigu plus de
deux tons majeurs.

Dans \struck{cette experience} \add{le cas présent} l'addition se fesant de part
et d'autre dela ligne nodale, celle-ci ne doit
pas se déplacer; et c'est aussi ceque montre l'experience

\page{545}

Une chose fort remarquable dans le phénomène dont je viens
de rendre compte, est qu'au moyen d'un poids ajouté, on
puisse obtenir d'une partie très petite comprise entre le côté
chargé et la ligne nodale parallelle la plus voisine de ce
côté, le même son que rend en même tems une partie \add{beaucoup plus grande}
prise sur la même plaque et comprise entre deux
lignes egalement parallelles au côté chargé.
\note{« beaucoup » est écrit en bout de ligne, « plus grande » dans l'interligne au-dessous, d'une plume plus fine. Sur cette page, « parallelle » et « nodale » (deux fois) et « transversale » portent sous la dernière lettre la petite croix qui annule un pluriel ; on donne le singulier.}

Ainsi, par exemple, j'ai pris une lame de verre
longue de 20 centimetres et large de 23 millimetres.
cette lame pesoit 4 gros 12 grains. La figure \struck{form} composée
d'une seule ligne nodale suivant la longueur, et de
\uncertain{2} suivant la largeur, s'\add{y} est formée très régulièrement, \struck{sur}
\struck{cette lame} c'est-àdire que l'espace entre les deux
lignes transversales, étoit exactement double de celui entre
l'une de ces lignes et l'extremité la plus voisine.
\note{Le chiffre lu « 2 » a une longue hampe descendante, et pourrait passer pour un 5 ; « les deux lignes transversales », deux lignes plus bas, appuie 2.}

Après avoir ajouté 1 gros 3 grains à une des lignes d'extremités
dela même lame, la ligne nodale la plus voisine parallelle
au côté chargé s'est formée à 12 millimetres de ce côté,
et la seconde ligne parallelle à 121 millimetres dela
premiere. Il restoit 61 millimetres environ entre cette
seconde ligne transversale et l'extremité opposée au
côté chargé, desorte que \struck{6 millimetres} apeuprès 6 millimetres
étoient occupés par la cire, la premiere ligne transversale et

\page{546}

et par la seconde des mêmes lignes bien plus large que
la premiere à cause de l'accumulation dela poussiere autour
du point d'application des doigts.

Je ne m'arrêterai pas à l'explication de cette experience.
je veux seulement faire remarquer que, dans le cas
présent, \struck{une portion} 12 millimetres pris entre le côté
chargé et la ligne parallelle la plus voisine, donnent
le même son que les 121 millimetres compris entre
les deux lignes transversales, et que pourtant une
de ces portions est dix fois plus \struck{p} étendue que
l'autre.
\note{Sur « son », la petite croix qui annule un pluriel ; de même plus bas sur « être » et « soumis ».}

J'ai reconnu qu'en général, on peut toujours en employant
\struck{un côté d'une plaque élastiq}
l'addition d'une matière pesante et non élastique; tirer
d'une portion \struck{comprise entre la ligne d'application}
\struck{et la ligne nodale parallelle la plus voisine} aussi petite
que l'on voudra \add{d'une plaque élastique,} le même son que d'une autre
portion de longueur donnée prise sur la même plaque
\note{« un côté d'une plaque élastiq », barré, est écrit dans l'interligne au-dessus de la fin de la première ligne ; « d'une plaque élastique, » dans l'interligne au-dessus de « le même son ». Un trait court sous « en employant », au niveau de la ligne, a été pris pour un trait de liaison et non pour une rature : sans ces mots la phrase n'aurait pas de verbe.}

L'explication physique de ce phénomène me par\uncertain{oît}
facile; en effet on conçoit qu'en fesant abstraction de
\uncertain{lagene}, qui \struck{tend à affoiblir le son} nuit à l'intensité
du son, les differens points dela plaque vibrante
doivent être absolument soumis aux mêmes conditions,
soit que le prolongement de cette plaque soit effectif
\note{La fin de « paroît » est couverte d'une tache. « lagene » : sans doute « la gêne », écrit d'un seul mot, sans accent visible.}

\page{547}

soit que l'addition d'une matière pesante \uncertain{et} capable
de faire equilibre àla portion dont la plaque devroit
être prolongée, tienne lieu de\struck{la} cette \struck{portion} prolongation.
Si ces deux cas \struck{peuvent} \add{doivent} être considérés comme entierement
équivalans, on peut concevoir la substitution dela
matière \struck{et} non élastique àla portion de surface
élastique qu'elle remplace, \struck{com} comme effectuée à une
distance quelconque dela ligne nodale; desorte que, comme
je l'ai dit plus haut, une partie aussi petite
que l'on voudra soit \uncertain{equivalante} à une autre
partie prise sur la même \struck{pl} surface et d'une
étendue donnée.
\note{Une tache d'encre couvre le milieu du mot lu « equivalante ». Le trait qui passe sur « portion », devant « prolongation », est léger ; on le tient pour une rature, les deux mots ne pouvant rester ensemble.}

\struck{Il est également facile de concevoir}

L'analyse montre également que l'addition dela cire
peut être faite à une distance quelconque de
l'origine. En effet, si on se rappelle que les
termes qui doivent être nuls \struck{sont nulls} aux limites
dela surface, sont tous \add{les deux} multipliés par un facteur
dependant de la rigidité naturelle dela matière
supposée elastique, on concevra que tout point dans
le quel ce facteur \struck{est} \add{sera} nul, \uncertain{satisfera} aux conditions
des extremités, independamment des valeurs des coordonnées
de la surface, \struck{c'est-àdire} \add{c'est-àdire} \struck{et parconséquent desorte}
que cette surface pourra alors se terminer partout où on voudra
\note{« satisfera » : la fin du mot est récrite à l'encre plus lourde, apparemment sur « satisfait » ; la lecture est offerte d'après « sera » ajouté au-dessus de « est ». « les deux » est écrit au-dessus de « tous ».}

\page{548}

\struck{Tout ce qui} Les applications que j'ai faites jusqu'ici, ne regardent
\add{généralement} \struck{qu} que les plaques rectangulaires, ou, \struck{lors} \add{si} elles comprennent aussi
d'autres figures, ce n'est \struck{qu'en considerant} \struck{qu'en considerant}
\struck{tant que ces dernieres peuvent être considerées comme parties}
\struck{integrantes d'une plaque rectangulaire en quelque sorte}
qu'accidentellement et sous certaines conditions particulières, dépendantes
del'état des extremités.

Il faudroit encore beaucoup de tems et de \struck{recherches} \add{travail} pour
déterminer egalement les \struck{figures et les sons} \add{manières de vibrer} dont sont
susceptibles des plaques de formes plus compliquées.
\struck{cependant avant de terminer ce me} je ne prétends
pas épuiser une matière si abondante; cependant, avant
de terminer ce memoire, je vais encore présenter
quelques recherches relatives au cas des plaques
circulaires, \struck{qui}
\note{« matière » et « memoire » portent sur la dernière lettre la petite croix qui annule un pluriel.}

§ 7. \emph{Recherche d'une intégrale particulière immédiatement applicable}
\emph{au cas d'une plaque elastique circulaire}
\note{Le titre est souligné. Le signe lu « § » a la forme d'un D pointé ; « d'une » est récrit sur un premier « D ».}

N.o \uncertain{21}. Il faut commencer par transformer l'équation différentielle
des surfaces élastiques vibrantes. En prenant donc le plan
des $\underline{x}$ et $\underline{y}$ parallelle au plan de la plaque en repos,
designant par $\underline{r}$ le rayon du cercle et par $\underline{\varphi}$ l'angle
qu'un rayon passant par le point auquel appartiennent
les coordonnées $\underline{x}$ et $\underline{y}$ \add{fait avec un des axes,} on pourra toujours faire
\note{Le numéro lu « 21 » : le premier chiffre a la boucle basse de ses 2 ; on ne peut exclure 31.}

\page{549}

\note{Un autre papier que la mise au net qui précède : une feuille de brouillon, écrite d'une main rapide, très raturée, en deux encres. La première ligne et les quatre suivantes sont d'une encre pâle ; le reste d'une encre noire. Dans la marge de gauche, le cachet de la Bibliothèque (« Bibliothèque impériale, MSS »), non transcrit. En haut à droite, à l'encre, le nombre 270.}

$z = \dots \{ e^{\frac{\omega x}{a}} - e^{\frac{(a-\omega)}{a}x} + (1$
\note{Ligne isolée en tête de la feuille, d'une encre pâle, interrompue après « (1 ». Elle n'a pas de rapport visible avec ce qui suit.}

$D$ Si après avoir effacé cequi se détruit on divise
par \add{\uncertain{$+\frac{dv}{dr}$}} la même équation se reduira à

$\frac{dv}{dr} + \frac{h}{k}\,v = 0 \dots (D)$
\note{La capitale lue D, dans la marge devant « Si » et dans le renvoi « (D) », a la forme d'un dôme posé sur un trait ; on pourrait y voir un $\Omega$. La lecture D s'accorde avec « $\frac{K}{CD}$ » plus bas, où la même lettre revient. L'addition au-dessus de « par la » est traversée d'un trait qui peut être une rature.}

Il est évident que l'on doit \struck{parvenir} obtenir, \uncertain{deux} \struck{deux}
\struck{On parviendroit} également \struck{aux} les equations (B) et (D)
en traitant la question d'une manière spéciale

Pour y parvenir il faut remarquer que dans l'hypothèse
présente le mouvement dela chaleur a lieu àla fois dans deux
directions differentes. \struck{suivant la premiere elle penetre dans une}
\struck{portion infiniment peu épaisse du cylindre, la seconde}
car elle penetre dans \struck{une} la tranche infiniment petite \struck{d'une}
\add{qui appartient à un disque d'une épaisseur infiniment petite aussi}
couche \add{\struck{parallele}} de ce cylindre, \struck{infiniment petite elle même} \add{\ill{}} \struck{de la \ill{} dela chaleur}
\note{Passage très surchargé. « qui appartient à un disque d'une épaisseur infi- » est écrit dans l'interligne au-dessus de « infiniment petite elle même », barré ; la fin du mot descend en ondulant le long du bord droit, et « aussi » est écrit au-dessus. Un petit mot abrégé, illisible, au-dessus de « même ».}

En appliquant ici l'analyse employée N.o \uncertain{19}
l'accroissement \struck{que} dela chaleur \struck{qui reçoit} pendant l'instant
$dt$ \struck{résultant} \add{est dû} \struck{du} \add{au} mouvement qui s'exécute dans la couche
comprise entre la surface dont le rayon est $r$ et \struck{en appliquant ici l'analyse}
\struck{employée dans cette \ill{} N.o \ill{}}
celle dont le rayon est $r+dr$, fournit dans \struck{l'équation} \add{(C)} les deux termes $\frac{K}{CD}\left(\frac{d^{2}v}{dr^{2}} + \frac{1}{r}\,\frac{dv}{dr}\right)$.
\note{« En appliquant ici l'analyse employée N.o 19 » est écrit dans l'interligne, au-dessus de « l'accroissement ». Devant « les deux termes », dans la marge, un mot barré suivi de « (C) », et au-dessus « \ill{} (C) » où le premier groupe est barré.}

\add{l'analyse employée N.o \uncertain{14} peut servir à prouver que}
\struck{On voit également}, le troisième terme $\frac{K}{CDrr}\,\frac{dv}{d\varphi}$
du second membre dela même équation (C) est dû
au mouvement \struck{qui} dela chaleur qui pendant l'instant
$dt$ s'exécute dans la tranche infiniment petite
comprise entre une section placée àla distance $r\varphi$
et une autre section placée àla distance $r\varphi + r\,d\varphi$

\struck{l'analyse employée N.o 14}
\note{Au-dessus de la fin de la ligne du troisième terme, un « 2 » isolé semble l'exposant oublié de $\frac{d^{2}v}{d\varphi^{2}}$ ; la fraction est donnée telle qu'elle est écrite. Dans la marge de gauche, en colonne, des essais de plume : un mot répété plusieurs fois, lu « tranche » par endroits, puis « on voit » barré et « que ». La page s'arrête sur la ligne barrée, un trait court au-dessous.}

\page{553}

\note{Deux petits feuillets de calculs collés sur une même feuille de montage, l'un en haut, l'autre en bas, chacun avec son nombre à l'encre en haut à droite : 273 sur celui du haut, 272 sur celui du bas. On les donne dans l'ordre où la vue les présente, de haut en bas.}

\[ \frac{dz}{dx} = \text{cos}\,\varphi\,\text{cos}\,\delta\,\frac{dz}{dr} - \frac{\text{sin.}\,\varphi\,\text{cos}\,\delta}{r}\,\frac{dz}{d\varphi} - \frac{\text{sin.}\,\delta\,\text{cos}\,\varphi}{r}\,\frac{dz}{d\delta} \]
\[ \frac{dz}{dy} = \text{sin}\,\varphi\,\text{sin.}\,\delta\,\frac{dz}{dr} + \frac{\text{cos}\,\varphi\,\text{cos.}\,\delta}{r}\,\frac{dz}{d\varphi} - \text{sin}\,\delta\,\text{sin.}\,\varphi\,\frac{dz}{dr} \]
\[ \frac{d^{2}z}{dx^{2}} = \frac{\text{sin}^{2}\varphi}{r^{\ill{}}}\,\text{cos.}\,\delta\,\frac{dz}{dr} \]
\note{Feuillet du haut (273), écrit sur son tiers supérieur seulement. Au dernier terme de la deuxième ligne le feuillet porte $\frac{dz}{dr}$ où la première ligne a $\frac{dz}{d\delta}$ ; donné tel quel. L'exposant de $r$ à la troisième ligne est couvert d'une tache.}

En prenant le rayon constant $r = a$
\[ \frac{dz}{dx} = -\,\frac{\text{sin}\,\varphi}{a}\,\frac{dz}{d\varphi} \]
\[ \frac{d\left(\frac{ddz}{dx^{2}} + \frac{ddz}{dy^{2}}\right)}{dx} = -\,\frac{\text{sin}\,\varphi}{a^{3}}\,\frac{d^{2}z}{d\varphi} \]
l'equation des extrêmités
\[ \left(\frac{d^{2}z}{dx^{2}} + \frac{d^{2}z}{dy^{2}}\right)\delta\frac{dz}{dx} - \frac{d\left(\frac{ddz}{dx^{2}} + \frac{ddz}{dy^{2}}\right)}{dx}\,\delta z = 0 \]
sera satisfaite par la condition $\text{sin}\,\varphi = 0$
qui donne en même tems $dx = a\,\text{sin}\,\varphi\,d\varphi = 0$
cecas est celui dans lequel on ne
considere que la circonference dela
plaque ronde.
\note{Feuillet du bas (272). Au dénominateur de la deuxième équation, un « $d\varphi$ » là où l'on attendrait $d\varphi^{2}$ ; à la première, « $dx$ » est écrit sur un premier « $da$ ». Donné tel qu'écrit.}

La condition qui vient du terme \uncertain{$\int\!\int$}
est $\frac{d^{2}z}{dx\,dy}\,\delta z = 0$ ainsi lorsque $\delta z$
n'est pas nul il faut qu'aux extremités
$\frac{d^{2}z}{dx\,dy} = 0$ ; $\int \frac{d^{2}z}{dx\,dy}\,dx$ \struck{$= \frac{dz}{dy} = A$}

\struck{$\int \frac{dz}{dy}\,dy = z = Ay + B$ ; $\frac{dz}{dy}\,dy = A\,dy$}
\note{Les deux dernières lignes du feuillet, de calcul, sont barrées ; la seconde est serrée au bas du papier.}

$\frac{d^{2}z}{dx\,dy} = 0$ donne $z = F(x) + f(y)$
c'est àdire que si les points de limites
satisfont aux \add{conditions des} extrêmités sans que $\delta z = 0$
$z$ sera \struck{une} \add{la somme de deux} fonctions séparées de $x$ et $y$
sauf les cas particuliers qui pourraient
rendre $\frac{d^{2}z}{dx\,dy}$ nul.

\struck{$\int \frac{d^{2}z}{dx\,dy}\,dy = \int A\,dy$} ; \struck{$z =$}

\struck{$\frac{dz}{dx} = Ay$}

\struck{$\int \frac{dz}{dx}\,dx = \int Ay\,dx + B$}

\struck{$z = Ayx + B$}
\note{Cette colonne occupe la moitié droite du feuillet du bas, à la hauteur de la fin de la colonne de gauche. « la somme de deux » est écrit au-dessus de « une », barré.}

\page{555}

\note{Un petit feuillet de calculs, sans phrase, collé au bas d'une feuille de montage ; en haut à droite, à l'encre, 274 ; le cachet de la Bibliothèque à droite du calcul. La lettre lue $\pi$ dans les arguments est un petit « n » surmonté d'un trait, suivi d'un point : « $\pi.nx$ ». Trois signes, dans l'avant-dernière ligne, sont couverts de taches ; ils sont marqués \ill{}.}

cas particuliers de l'intégrale (E) ; \uncertain{$n = -1$} ; $b = B = 0$
\note{Après « $-1$ », un trait descendant bouclé, qui peut être une virgule, un point-virgule ou un chiffre ; la barre sur le second B peut être un accent.}

\[ z = \dots\ \text{sin.}\frac{\pi nx}{A}\,\text{sin.}\frac{\pi my}{A}\left\{a'\left(\text{cos}\frac{\pi nx}{A} - 2\,\text{cos.}^{3}\frac{\pi nx}{A}\right)\ill{}\,A'\left(\text{cos}\frac{\pi my}{A} - 2\,\text{cos.}^{3}\frac{\pi my}{A}\right)\right\} \]

\[ \frac{dz}{dx} = \dots\ \frac{\pi.n}{A}.\,\text{cos}\frac{\pi.nx}{A}\,\text{sin}\frac{\pi.my}{A}\left\{a'\left(\text{cos.}\frac{\pi.nx}{A} - 2\,\text{cos.}^{3}\frac{\pi.nx}{A}\right)A'\left(\text{cos}\frac{\pi.my}{A} - 2\,\text{cos.}^{3}\frac{\pi my}{A}\right)\right\} \]
\[ \dots -\frac{\pi n}{A}\,\text{sin}\frac{\pi.nx}{A}\,\text{sin}.\pi.my\left\{\left(\text{sin}\frac{\pi.nx}{A} - 6\,\text{cos}^{2}\frac{\pi.nx}{A}\right)A'\left(\text{cos.}\frac{\pi.my}{A} - 2\,\text{cos.}^{3}\frac{\pi my}{A}\right)\right\} \]
\note{Dans la première ligne de $\frac{dz}{dx}$, « $\text{cos}\frac{\pi.nx}{A}$ » est récrit sur un autre mot. Dans la seconde, « $\text{sin}\,\pi.nx$ » est écrit au-dessus de « $6\,\text{cos}^{2}$ » ; « $\text{sin}.\pi.my$ » est sans dénominateur. Les deux lignes sont données telles qu'elles sont écrites, sans les rapprocher du calcul de la dérivée.}

\[ \frac{ddz}{dx^{2}} = \dots -\frac{\pi^{2}n^{2}}{A^{2}}\,\text{sin}\frac{\pi.nx}{A}\,\text{sin}\frac{\pi.my}{A}\left\{a'\left(\text{cos}\frac{\pi.nx}{A} - 2\,\text{cos}^{3}\frac{\pi.nx}{A}\right)A'\left(\text{cos.}\frac{\pi.my}{A} - 2\,\text{cos.}^{3}\frac{\pi.my}{A}\right)\right. \]
\[ -\,2\,\frac{\pi^{2}n^{2}}{A^{2}}\,\text{sin}\frac{\pi.nx}{A}\,\text{sin}\frac{\pi.my}{A}\left\{\text{cos}\frac{\pi.nx}{A} - 6\,\text{cos.}^{3}\frac{\pi.nx}{A}\right\} \]
\[ \ill{}\ \frac{\pi n^{2}}{A}\,\text{sin}\frac{\pi.nx}{A}.\,\text{sin}\frac{\pi.my}{A}\left\{\text{cos.}\frac{\pi.nx}{A}\ \ill{}\ 12\,\struck{\text{sin}}^{2}\frac{\pi.nx}{A}\,\text{cos.}\frac{\pi.nx}{A}\ \ill{}\ 6\,\text{cos}^{3}\frac{\pi x n}{A}\right\} \]
\[ = \dots -\frac{\pi^{2}n^{2}}{A^{2}}.\,\text{sin}\frac{\pi.nx}{A}\,\text{sin}\frac{\pi.my}{A}\left\{16.\,a'\left(\text{cos.}\frac{\pi.nx}{A} - 2\,\text{cos.}^{3}\frac{\pi.nx}{A}\right)A'\left(\text{cos.}\frac{\pi.my}{A} - 2\,\text{cos.}^{3}\frac{\pi.my}{A}\right)\right\} \]
\note{La première ligne de $\frac{ddz}{dx^{2}}$ n'a pas d'accolade fermante sur le feuillet. Dans l'avant-dernière, les trois signes couverts de taches séparent les termes ; au second, un mot barré devant l'exposant 2, et « cos. » repassé à l'encre lourde ; le dernier argument est écrit « $\pi x n$ ». Le feuillet s'arrête après la ligne en « 16. ».}

\page{557}

\note{Un feuillet de calculs mis au net, d'une écriture régulière, presque sans rature. En tête, au milieu, entre parenthèses, « (138) » barré d'un trait horizontal ; en haut à droite, à l'encre, 275. Les parenthèses qui entourent les coefficients différentiels sont les siennes.}

$x = r\,\text{cos}\,\varphi$, $y = r\,\text{sin.}\,\varphi$ et a cause de $\frac{d\varphi}{dx} = -\,\frac{\text{sin}\,\varphi}{r}$, $\frac{d\varphi}{dy} = \frac{\text{cos}\,\varphi}{r}$

$\frac{dr}{dx} = \text{cos}\,\varphi$, ; $\frac{dr}{dy} = \text{sin}\,\varphi$

on trouvera:
\[ \frac{dz}{dx} = \text{cos.}\,\varphi.\left(\frac{dz}{dr}\right) - \frac{\text{sin.}\,\varphi}{r}\left(\frac{dz}{d\varphi}\right) \]
\[ \frac{dz}{dy} = \text{sin.}\,\varphi\left(\frac{dz}{dr}\right) + \frac{\text{cos}\,\varphi}{r}\left(\frac{dz}{d\varphi}\right) \]
\[ \frac{ddz}{dx^{2}} = \frac{\text{sin}^{2}\varphi}{r}\left(\frac{dz}{dr}\right) + \frac{2\,\text{sin.}\,\varphi\,\text{cos}\,\varphi}{rr}\left(\frac{dz}{d\varphi}\right) + \text{cos}^{2}\varphi\left(\frac{ddz}{dr^{2}}\right) - \frac{2\,\text{sin.}\,\varphi\,\text{cos}\,\varphi}{r}\left(\frac{ddz}{dr\,d\varphi}\right) \]
\[ + \frac{\text{sin}^{2}\varphi}{r^{2}}\left(\frac{ddz}{d\varphi^{2}}\right) \]
\[ \frac{ddz}{dy^{2}} = \frac{\text{cos}^{2}\varphi}{r}\left(\frac{dz}{dr}\right) - \frac{2\,\text{sin.}\,\varphi\,\text{cos.}\,\varphi}{rr}\left(\frac{dz}{d\varphi}\right) + \text{sin.}^{2}\varphi\left(\frac{ddz}{dr^{2}}\right) + \frac{2\,\text{sin.}\,\varphi\,\text{cos.}\,\varphi}{r}\,\frac{ddz}{dr\,d\varphi} + \frac{\text{cos}^{2}\varphi}{r^{2}}\left(\frac{ddz}{d\varphi^{2}}\right) \]
\note{Dans $\frac{ddz}{dy^{2}}$, le facteur $\frac{ddz}{dr\,d\varphi}$ du quatrième terme est écrit au-dessus de la ligne, entre la fraction et le signe qui suit.}
\[ \frac{ddz}{dx^{2}} + \frac{ddz}{dy^{2}} = \frac{1}{r}\left(\frac{dz}{dr}\right) + \left(\frac{ddz}{dr^{2}}\right) + \frac{1}{rr}\left(\frac{ddz}{d\varphi^{2}}\right) \]
\[ \frac{d\left\{\frac{ddz}{dx^{2}} + \frac{ddz}{dy^{2}}\right\}}{dx} = -\,\frac{\text{cos.}\,\varphi}{r^{2}}.\frac{dz}{dr} + \frac{\text{cos.}\,\varphi}{r}.\frac{ddz}{dr^{2}} - \frac{\text{sin.}\,\varphi}{r^{2}}.\frac{ddz}{dr\,d\varphi} + \text{cos.}\,\varphi.\frac{d^{3}z}{dr^{3}} \]
\[ - \frac{\text{sin.}\,\varphi}{r}\,\frac{d^{3}z}{dr^{2}d\varphi} + \frac{\text{cos.}\,\varphi}{r^{2}}.\frac{d^{3}z}{dr\,d\varphi^{2}} - \frac{\text{sin.}\,\varphi}{r^{3}}.\frac{d^{3}z}{d\varphi^{3}} - \frac{2\,\text{cos.}\,\varphi}{r^{3}}.\frac{ddz}{d\varphi^{2}} \]
\note{Au dernier terme de la première ligne, l'exposant du dénominateur, $dr^{3}$, est récrit et taché.}
\[ \frac{dd\left\{\frac{ddz}{dx^{2}} + \frac{ddz}{dy^{2}}\right\}}{dx^{2}} = -\left\{\frac{\text{sin}^{2}\varphi}{r^{3}} - \frac{2\,\text{cos}^{2}\varphi}{r^{3}}\right\}.\frac{dz}{dr} + \left\{\frac{\text{sin}^{2}\varphi}{r^{2}} - \frac{2\,\text{cos}^{2}\varphi}{r^{2}}\right\}\frac{ddz}{dr^{2}} + \frac{4\,\text{sin.}\,\varphi\,\text{cos.}\,\varphi}{r^{3}}.\frac{ddz}{dr\,d\varphi} \]
\[ + \left\{\frac{\text{sin}^{2}\varphi}{r} + \frac{\text{cos}^{2}\varphi}{r}\right\}\frac{d^{3}z}{dr^{3}} - \left\{\frac{2\,\text{sin}^{2}\varphi}{r^{4}} - \frac{6\,\text{cos}^{2}\varphi}{r^{4}}\right\}\frac{ddz}{d\varphi^{2}} + \frac{6\,\text{sin}\,\varphi\,\text{cos.}\,\varphi}{r^{4}}.\frac{d^{3}z}{d\varphi^{3}} \]
\[ + \left\{\frac{2\,\text{sin}^{2}\varphi}{r^{3}} - \frac{4\,\text{cos}^{2}\varphi}{r^{3}}\right\}.\frac{d^{3}z}{d\varphi^{2}dr} + \text{cos}^{2}\varphi.\frac{d^{4}z}{dr^{4}} - 2\,\frac{\text{sin.}\,\varphi\,\text{cos}\,\varphi}{r}.\frac{d^{4}z}{dr^{3}d\varphi} \]
\[ + \left\{\frac{\text{sin}^{2}\varphi}{r^{2}} + \frac{\text{cos}^{2}\varphi}{r^{2}}\right\}.\frac{d^{4}z}{d\varphi^{2}dr^{2}} - 2\,\frac{\text{sin.}\,\varphi\,\text{cos.}\,\varphi}{r^{3}}.\frac{d^{4}z}{d\varphi^{3}dr} + \frac{\text{sin.}^{2}\varphi}{r^{4}}.\frac{d^{4}z}{\uncertain{d\varphi^{4}}} \]
\note{Le « 4 » du coefficient $\frac{4\,\text{sin.}\,\varphi\,\text{cos.}\,\varphi}{r^{3}}$ et celui de $r^{4}$ ont la forme de « b » que prend son 4. Le dernier dénominateur touche le bord du feuillet.}
\[ \frac{d\left\{\frac{ddz}{dx^{2}} + \frac{ddz}{dy^{2}}\right\}}{dy} = -\,\frac{\text{sin.}\,\varphi}{r^{2}}.\frac{dz}{dr} + \frac{\text{sin.}\,\varphi}{r}.\frac{ddz}{dr^{2}} + \frac{\text{cos.}\,\varphi}{r^{2}}.\frac{ddz}{dr.d\varphi} + \text{sin}\,\varphi.\frac{d^{3}z}{dr^{3}} \]
\[ + \frac{\text{cos.}\,\varphi}{r}.\frac{d^{3}z}{dr^{2}d\varphi} + \frac{\text{sin.}\,\varphi}{r^{2}}.\frac{d^{3}z}{dr\,d\varphi^{2}} + \frac{\text{cos.}\,\varphi}{r^{3}}.\frac{d^{3}z}{d\varphi^{3}} - 2\,\frac{\text{sin.}\,\varphi}{r^{3}}.\frac{ddz}{d\varphi^{2}} \]
\note{Le bas du feuillet est blanc. Sur toute la page, les dérivées sont données avec les signes et les coefficients du feuillet, sans vérification.}

\page{560}

\note{Le dos du feuillet 275 (vues 557 et 559) : en tête, au milieu, « (139) », non barré ; pas de nombre en haut à droite. La page suit directement le calcul de la page (138).}

\[ \frac{dd\left\{\frac{ddz}{dx^{2}} + \frac{ddz}{dy^{2}}\right\}}{dy^{2}} = -\left\{\frac{\text{cos}^{2}\varphi}{r^{3}} - \frac{2\,\text{sin.}^{2}\varphi}{r^{3}}\right\}.\frac{dz}{dr} + \left\{\frac{\text{cos}^{2}\varphi}{r^{2}} - \frac{2\,\text{sin.}^{2}\varphi}{r^{2}}\right\}.\frac{ddz}{dr^{2}} - \frac{4\,\text{cos.}\,\varphi\,\text{sin.}\,\varphi}{r^{3}}.\frac{ddz}{dr\,d\varphi} \]
\[ + \left\{\frac{\text{cos}^{2}\varphi}{r} + \frac{\text{sin}^{2}\varphi}{r}\right\}\frac{d^{3}z}{dr^{3}} - \left\{\frac{2\,\text{cos}^{2}\varphi}{r^{4}} - \frac{6\,\text{sin}^{2}\varphi}{r^{4}}\right\}\frac{ddz}{d\varphi^{2}} - \frac{6\,\text{sin.}\,\varphi\,\text{cos.}\,\varphi}{r^{4}}.\frac{d^{3}z}{d\varphi^{3}} \]
\[ + \left\{\frac{2\,\text{cos}^{2}\varphi}{r^{3}} - \frac{4\,\text{sin}^{2}\varphi}{r^{3}}\right\}\frac{d^{3}z}{d\varphi^{2}dr} + \text{sin}^{2}\varphi.\frac{d^{4}z}{dr^{4}} + \frac{2\,\text{sin.}\,\varphi\,\text{cos.}\,\varphi}{r}.\frac{d^{4}z}{dr^{3}d\varphi} \]
\[ + \left\{\frac{\text{cos}^{2}\varphi}{r^{2}} + \frac{\text{sin}^{2}\varphi}{r^{2}}\right\}\frac{d^{4}z}{dr^{2}d\varphi^{2}} + \frac{2\,\text{sin.}\,\varphi\,\text{cos.}\,\varphi}{r^{3}}.\frac{d^{4}z}{dr\,d\varphi^{3}} + \frac{\text{cos}^{2}\varphi}{r^{4}}.\frac{d^{4}z}{d\varphi^{4}} \]

\[ \frac{d^{4}z}{dx^{4}} + 2\,\frac{d^{4}z}{dx^{2}dy^{2}} + \frac{d^{4}z}{dy^{4}} = \frac{1}{r^{3}}.\frac{dz}{dr} - \frac{1}{r^{2}}.\frac{ddz}{dr^{2}} + \frac{4}{r^{4}}.\frac{ddz}{d\varphi^{2}} + \frac{2}{r}.\frac{d^{3}z}{dr^{3}} \]
\[ - \frac{2}{r^{3}}.\frac{d^{3}z}{d\varphi^{2}dr} + \frac{d^{4}z}{dr^{4}} + \frac{2}{r^{2}}.\frac{d^{4}z}{dr^{2}d\varphi^{2}} + \frac{1}{r^{4}}.\frac{d^{4}z}{d\varphi^{4}} \]
\note{Le $z$ du premier terme et celui de $\frac{d^{3}z}{dr^{3}}$ sont récrits et tachés.}

En fesant donc \struck{k} $b\tau^{2} = \frac{e^{4}}{a^{4}}$ l'equation différentielle (B) (V. N.o 2)
sera de la forme
\note{La lettre lue $\tau$, l'épaisseur, est le trait haut coiffé d'un crochet déjà relevé dans le mémoire sur les plaques (vues 402–413) ; ici elle ressemble à un $r$ ou à un $\Gamma$. La lettre lue $e$ au numérateur pourrait être un $c$.}

\struck{$\frac{a^{4}}{e^{4}}\ \ill{}\ \frac{1}{r^{3}}\frac{dz}{dr}\ \ill{}\ \frac{\ill{}}{r^{2}}$}
\[ \frac{a^{4}}{e^{4}}.z = \struck{m}\left\{\begin{array}{l} \frac{1}{r^{3}}.\frac{dz}{dr} - \frac{1}{r^{2}}\,\frac{ddz}{dr^{2}} + \frac{2}{r}\,\frac{d^{3}z}{dr^{3}} + \frac{d^{4}z}{dr^{4}} \\ - \frac{2}{r^{3}}.\frac{d^{3}z}{d\varphi^{2}dr} + \frac{2}{r^{2}}.\frac{d^{4}z}{dr^{2}d\varphi^{2}} + \frac{4}{r^{4}}\,\frac{ddz}{d\varphi^{2}} + \frac{1}{r^{4}}\,\frac{d^{4}z}{d\varphi^{4}} \end{array}\right. \]
\note{Un premier essai de la même équation, en tête de la ligne, est barré et en partie noirci. Après le signe $=$, une lettre barrée, lue « m ». L'accolade est ouverte à gauche des deux lignes et ne se ferme pas ; le signe devant $\frac{1}{r^{2}}$ est récrit sur une tache.}

cette equation ne contient\struck{t} que des differentielles partielles d'un ordre
pair de l'angle $\varphi$, \struck{$u$ étant une fonction de $r$ seule} on peut
faire $z = u\,\text{sin}(\beta\varphi + B)$, $u$ étant une fonction de $r$ seul.

Si, après avoir substitué cette valeur de $\underline{z}$ dans l'équation
précédente, on la divise par $\text{sin}(\beta\varphi + B)$, elle deviendra
\[ \frac{a^{4}}{e^{4}}\,u = \frac{1 + 2\beta\beta}{r^{3}}.\frac{du}{dr} - \frac{1 + 2\beta\beta}{r^{2}}.\frac{ddu}{dr^{2}} + \frac{2}{r}.\frac{d^{3}u}{dr^{3}} + \frac{d^{4}u}{dr^{4}} - \frac{4\beta\beta u}{r^{4}} + \frac{\beta^{4}u}{r^{4}} \dots (\text{B}') \]
\note{Le renvoi final est coupé au bord : « (B » et un trait, lu (B$'$). Le $\beta$ de ces lignes est une lettre bouclée, à hampe montante.}

Apresent on peut représenter par \struck{s} une serie formée des
puissances ascendantes de $r$, et supposer $u = r^{\beta}s$, \struck{et} \add{cette supposition donnera} \struck{par conséquent}
\note{« Apresent » : ainsi, en un mot. Après « par », une lettre noircie. « cette supposition donnera » est écrit au-dessus de « et par conséquent », barré.}
\[ \frac{du}{dr} = \beta r^{\beta-1}s + r^{\beta}\,\frac{ds}{dr} \]
\[ \frac{ddu}{dr^{2}} = \beta(\beta-1)r^{\beta-2}s + 2\beta r^{\beta-1}.\frac{ds}{dr} + r^{\beta}\,\frac{dds}{dr^{2}} \]
\[ \frac{d^{3}u}{dr^{3}} = \beta(\beta-1)(\beta-2)r^{\beta-3}s + 3\beta(\beta-1)r^{\beta-2}\,\frac{ds}{dr} + 3\beta.r^{\beta-1}\,\frac{dds}{dr^{2}} + r^{\beta}\,\frac{d^{3}s}{dr^{3}} \]
\note{La page s'arrête sur cette ligne ; le bas du feuillet est blanc.}

\end{document}
